\documentclass[11pt,a4paper]{article}
\usepackage[margin=27mm]{geometry}
\usepackage{amsmath,amsthm,amssymb,fontspec,unicode-math,hyperref}
\setmainfont{STIXGeneral}
\setmathfont{STIX Two Math}
\setmonofont{Menlo}
\hypersetup{hidelinks}
\newtheorem{theorem}{Theorem}
\newtheorem{lemma}[theorem]{Lemma}
\newtheorem{corollary}[theorem]{Corollary}
\newtheorem{proposition}[theorem]{Proposition}
\theoremstyle{remark}\newtheorem{remark}[theorem]{Remark}
\newcommand{\gen}[1]{\langle #1\rangle}
\newcommand{\normalgen}[1]{\langle\!\langle #1\rangle\!\rangle}
\newcommand{\F}{\mathbb F}
\newcommand{\Z}{\mathbb Z}
\DeclareMathOperator{\kerop}{ker}
\title{Recognizing presentations by normal supported kernels}
\author{Human Compression Project}
\date{11 October 2026}
\begin{document}
\maketitle
\begin{abstract}
The construction of a compatible product and the recognition of a proposed
presentation are different tasks. We give an order-free recognition criterion:
normal local kernels, separately detected by coordinates, account for the whole
kernel of a base map. This proves an actual isomorphism without first proving
that the candidate group is finite. It gives a common proof of the presentation
steps in the existing three-generator hypertope and square-section constructions.
The local normality calculations remain explicit. The general injectivity
criterion, direct-product isomorphism and quotient transport have been checked in Lean; the source-specific
presentation and geometric hypotheses have written proofs here.
\end{abstract}

\section{The recognition criterion}
All homomorphisms in this paper are group homomorphisms. A product of subgroups
means a product of sets until it has been shown to be a subgroup. No group is
assumed finite in this section.

\begin{theorem}[Normal kernel recognition]\label{thm:recognition}
Let $\pi:P\to B$ and $\rho_i:P\to E_i$, for $i$ in a finite set $J$, be
homomorphisms. Suppose $K_i\unlhd P$ satisfy
\begin{align}
 \ker\pi&=\gen{K_i:i\in J},\label{eq:basekernel}\\
 \rho_i(K_j)&=1\quad(i\ne j),\label{eq:offdiag}\\
 \ker\rho_i\cap K_i&=1.\label{eq:faithful}
\end{align}
Then the homomorphism
\[
 P\longrightarrow B\times\prod_{i\in J}E_i,
 \qquad x\longmapsto (\pi(x),(\rho_i(x))_i)
\]
is injective. Moreover, the kernel in \eqref{eq:basekernel} is the internal
direct product of the $K_i$. For injectivity alone it suffices to replace
\eqref{eq:basekernel} by $\ker\pi\leq\gen{K_i:i\in J}$.
\end{theorem}
\begin{proof}
For $i\ne j$, an element of $K_i\cap K_j$ is killed by $\rho_i$ by
\eqref{eq:offdiag}, and is therefore trivial by \eqref{eq:faithful}.
Normality gives $[K_i,K_j]\leq K_i\cap K_j$, so distinct factors commute.
Their finite set product is consequently their generated subgroup. If
$\prod_i k_i=1$ with $k_i\in K_i$, applying $\rho_j$ gives
$\rho_j(k_j)=1$, whence $k_j=1$. Thus multiplication identifies the external
product with the generated subgroup. Finally, an element in the kernel of the
displayed homomorphism belongs to $\ker\pi$, hence has such a product
expression. Every coordinate vanishes, so every factor is trivial.
\end{proof}

\begin{corollary}[Recognition of a quotient]\label{cor:quotient}
Let $f:P\twoheadrightarrow H$ be a homomorphism, and let $H$ have a base
homomorphism $p:H\to B$ and coordinates $c_i:H\to E_i$. If
Theorem~\ref{thm:recognition} applies to $p f$, $c_i f$ and normal subgroups
$K_i\leq P$, then $f$ is an isomorphism.
\end{corollary}
\begin{proof}
The coordinate map factors through $f$, so injectivity of the coordinate map
implies injectivity of $f$. Surjectivity is the stated hypothesis.
\end{proof}

Normality is substantial. Let $P=D_8=\gen{u,v\mid u^2=v^2=(uv)^4=1}$,
let $B=1$, and take $K_1=\gen u$, $K_2=\gen v$. Its abelianization
$P\to C_2^2$ has coordinates faithfully detecting the two factors and killing
the other one. The factors generate $P$, but the map has the nontrivial kernel
$\gen{(uv)^2}$. Here neither factor is normal. Thus a correct local model and
coordinate detection do not themselves eliminate mixed relations.

\section{Supported extensions and complete presentations}
We use the already proved supported-kernel theorem [UA, UA.SK], with the
following full interface. Let $B=\gen{b_t:t\in I}$ have the intersection
property. For each member $S$ of a finite antichain $\mathcal S$ of subsets
of $I$, there is a deletion $p_S:B\to B_S$, fixing the $S$-markers and
killing the others, and a marked epimorphism $q_S:E_S\to B_S$ that is
injective on every proper marked subgroup of $E_S$. Set $N_S=\ker q_S$.
The compatible product
\begin{equation}
 H=\{(b,(e_S)_S):q_S(e_S)=p_S(b)\text{ for every }S\}
 \leq B\times\prod_S E_S                                      \label{eq:fiber}
\end{equation}
is generated by its prescribed marker tuples. Its base kernel is
$\prod_S N_S$, and its marked subgroup on $L\subseteq I$ has kernel
$\prod_{S\subseteq L}N_S$. Its marking has the intersection property.
The antichain assumption makes each $H_S\to E_S$ a marked isomorphism.
These statements include generation of the \emph{whole} group in
\eqref{eq:fiber}; they are not assumptions that a subgroup has the ambient
product's order.

\begin{theorem}[Presentation recognition for an antichain]\label{thm:presentation}
In the preceding setting, let $P=\gen{r_t:t\in I}$ be given by a
presentation. Suppose:
\begin{enumerate}
\item its relators hold on the marker tuples in \eqref{eq:fiber};
\item the local relations of every $E_S$ hold on $(r_t)_{t\in S}$;
\item the image $K_S$ of $N_S$ under $E_S\to P$ is normal in $P$;
\item adjoining the relations $K_S=1$ for all $S$ gives exactly the marked
presentation of $B$.
\end{enumerate}
Then the marker map $P\to H$ is an isomorphism. In particular the $K_S$
form an internal direct product, all marked intersections are exact, and if
the input groups are finite then
\begin{equation}
 |P_L|=|B_L|\prod_{S\subseteq L}|N_S|.                         \label{eq:orders}
\end{equation}
The additional deletion conclusion of UA.SK.Delete transfers to $P$ under
its hypotheses. The same holds for the geometric conclusion of
UA.SK.Geometry under its hypotheses.
\end{theorem}
\begin{proof}
The first assumption and the generation assertion of UA.SK give a marked
surjection $f:P\twoheadrightarrow H$. The second gives $E_S\to P_S$;
its composite with $f$ and projection to $E_S$ is the identity on every
local marker, hence is the identity homomorphism. Thus each local map is
injective, without an order comparison. Under $f$, the image of $N_S$ is
its pure $S$-coordinate: the antichain makes every other intersection of
supports proper, where $q_T$ is injective. Consequently the coordinate maps
kill $K_S$ outside its own coordinate and are injective on $K_S$ in that
coordinate.

Normality makes $K=\gen{K_S:S\in\mathcal S}$ a normal subgroup. The last
assumption identifies the marked map $P/K\to B$ with an isomorphism, so
the base map $P\to B$ has kernel exactly $K$. Corollary~\ref{cor:quotient}
now proves that $f$ is an isomorphism. The remaining conclusions follow
from UA.SK and the explicitly indicated additional clauses.
\end{proof}

The last assumption is a presentation calculation. It does not assume the
isomorphism $P\cong H$ which is being proved. In practice it says that, after
killing the displayed local kernel words, the remaining relations are the
known base presentation. A pair of marked maps which are inverse on generators
also proves this condition. Normality likewise has a concrete local test: the
markers within $S$ already normalize $K_S$, and one must check the remaining
markers. It is enough that those remaining markers centralize a generating
set of $K_S$. Merely writing down the desired quotient does not supply this
check.

\section{The three-generator hypertope presentation}
We retain the proved affine square model [UA, UA.Affine, UA.Square and
UA.SquareDecomposition]. In coordinates, with the rightmost transformation
acting first, it is
\[
 a(x,y)=(-x,y),\qquad b(x,y)=(y,x),\qquad c(x,y)=(x,1-y).
\]
These generate $W=\Z^2\rtimes D_8$ with complete presentation
\begin{equation}
 \gen{a,b,c\mid a^2=b^2=c^2=(ab)^4=(bc)^4=(ac)^2=1}.           \label{eq:square}
\end{equation}
For $\delta\geq0$, define
\[
 Q_\delta=\begin{cases}
 W/q\Z^2,&\delta\text{ even},\quad q=2^{(\delta+2)/2},\\
 W/q\gen{(1,1),(1,-1)},&\delta\text{ odd},\quad q=2^{(\delta+1)/2}.
 \end{cases}
\]
It has order $2^{\delta+5}$, pair orders $(4,4,2)$ and all Boolean deletions.
To identify the parity relation exactly, use $[x,y]=x^{-1}y^{-1}xy$.
Direct affine multiplication gives translations
\[
 [(ab)^2,c]=(0,-2),\qquad [(ab)^2,(bc)^2]=(-2,-2).
\]
The normal closure of a translation in $W$ is the lattice spanned by its
signed-permutation orbit. Hence, respectively, adding
\begin{equation}
 [(ab)^2,c]^{2^{\delta/2}}=1\quad\hbox{or}\quad
 [(ab)^2,(bc)^2]^{2^{(\delta-1)/2}}=1                          \label{eq:parity}
\end{equation}
to \eqref{eq:square} gives exactly $Q_\delta$.

Take $n\geq10$, $s,t\geq2$, $\ell\geq1$, and
$\delta=n-s-t-\ell\geq0$. Let $P$ have the three involutory markers,
the pair-power relations
\[
 (ab)^{2^s}=(bc)^{2^t}=(ac)^{2^\ell}=1,
\]
the three relations
\begin{equation}
 [(ab)^4,c]=[a,(bc)^4]=[(ac)^2,b]=1,                           \label{eq:normal}
\end{equation}
and the appropriate parity relation \eqref{eq:parity}. This is the defining
presentation of [HFL19, Theorem3.2], with renamed markers.

\begin{corollary}\label{cor:hypertope}
The presentation $P$ is the supported extension of $Q_\delta$ by the three
marked dihedral quotients of orders $2^{s+1},2^{t+1},2^{\ell+1}$ onto
its pair subgroups. In particular it has order $2^n$, the prescribed pair
orders, and a regular rank-three hypertope with full type-preserving group
$P$. The subgroups
\[
 K_{ab}=\gen{(ab)^4},\qquad K_{bc}=\gen{(bc)^4},\qquad
 K_{ac}=\gen{(ac)^2}
\]
form its direct-product base kernel, of type
$C_{2^{s-2}}\times C_{2^{t-2}}\times C_{2^{\ell-1}}$.
\end{corollary}
\begin{proof}
The three pair supports are an antichain. Their proper marked subgroups
are the trivial group and the reflection subgroups, so the maps to the
base are injective on every proper marked subgroup. UA.SK therefore
constructs $H$.

All relators of $P$ hold in $H$. Each pair-power kernel word is supported
only in its own coordinate, where the third marker is the identity, proving
\eqref{eq:normal}. The two commutator words in \eqref{eq:parity} vanish in
each local pair coordinate: one of the relevant markers is deleted, and
direct substitution gives the identity. Their base values are the specified
translations. Each local dihedral presentation holds in $P$. Within its
pair, each marker inverts the displayed kernel generator; the third marker
centralizes it by \eqref{eq:normal}. Thus all three $K$'s are normal.
Killing them reduces the presentation precisely to \eqref{eq:square} and
\eqref{eq:parity}. Theorem~\ref{thm:presentation} applies.

Its order formula gives
$2^{\delta+5}2^{s-2}2^{t-2}2^{\ell-1}=2^n$ and its local recognition gives
the exact pair orders. The square base is a regular string polytope; its
maximal-parabolic cosets have the needed flag property by the retained
polytope correspondence. The dihedral local markers are involutions, so
UA.SK.Geometry supplies the claimed hypertope and its full group. This
last use is explicit: intersection property alone does not supply the flag
property of a general coset geometry.
\end{proof}

The quotient description is consequently
\[
 P\cong
 (C_{2^{s-2}}\times C_{2^{t-2}}\times C_{2^{\ell-1}})\mathbin{.}
 \begin{cases}
 (D_{2^{(\delta+4)/2}}\times D_{2^{(\delta+4)/2}})\rtimes C_2,
       &\delta\text{ even},\\
 (D_{2^{(\delta+3)/2}}\times D_{2^{(\delta+3)/2}})\rtimes C_2^2,
       &\delta\text{ odd}.
 \end{cases}
\]
Here the extension dot makes no splitting assertion. These are the corrected
parity rows already recorded in the accepted source contract; the source
interchanges them and repeats the first pair word when naming its second
kernel. The present argument preserves both corrections. The decompositions
of the square quotient are the previously proved UA.SquareDecomposition,
not new group classifications.

\section{Square sections along a string}
Let $d\geq3$ and $n_3,\ldots,n_d\geq5$. Write
$S_i=\{i,i+1,i+2\}$ for $0\leq i\leq d-3$, and let $E_i=H_{n_{i+3}}$
be the square group of UA.Square. Its base quotient is $H_5$, and its kernel
is $\gamma_3(E_i)$, generated by its two doubled coordinate translations.
Those translations need not generate independent cyclic factors: in $H_6$
they coincide. Their full subgroup, rather than a presumed product of two
cyclic orders, is used throughout.

The common base $B$ is the class-two path group with $d$ binary vertex
coordinates and $d-1$ binary edge coordinates. Its multiplication is
\[
 (v,z)(w,t)=(v+w,z+t+\beta(v,w)),\qquad
 \beta_i(v,w)=v_{i+1}w_i.
\]
The full coordinate model and presentation are proved in UA.Forest: its order
is $2^{2d-1}$, its marked subgroups retain exactly the internal vertex and
edge coordinates, and it has every Boolean deletion. Each $B_{S_i}=H_5$.
The quotient $E_i\to H_5$ is faithful on every proper marked subgroup, since
all pair orders remain $4,4,2$. Thus UA.SK gives a group $H$ with precisely
these square sections and order
\begin{equation}
 |H|=2^{2d-1+\sum_{j=3}^d(n_j-5)}
     =2^{\sum_{j=3}^d n_j-3(d-3)}.                            \label{eq:pathorder}
\end{equation}
This also identifies $H$ with the particular iterated mix constructed in
[CFHS25, Theorem5.1], not merely with a group of the same order and local
type. Indeed, the projection $H\to\prod_i E_i$ is injective. If every local
coordinate is trivial, then every triple deletion of its base coordinate is
trivial. Every vertex and edge of the path lies in a consecutive triple, so
the displayed coordinates of $B$ force that base element to be trivial.
The projected image is generated by the marker tuples in the local factors.
The source's definition (8) and its induction in Theorem5.1 produce exactly
these tuples, with only a reassociation of the direct product. Hence they
give the same marked group.

Let $P_0$ have the string Coxeter relations of type $\{4,\ldots,4\}$ and
the complete local lattice relations of each $E_i$. Let $P$ be its quotient
by
\begin{equation}
 [(r_i r_{i+1})^2,(r_{i+2}r_{i+3})^2]=1
       \qquad(0\leq i\leq d-4).                             \label{eq:four}
\end{equation}

\begin{lemma}[The local normalization calculation]\label{lem:four}
If involutions $a,b,c,d$ satisfy $[a,c]=[a,d]=[b,d]=1$ and
$[(ab)^2,(cd)^2]=1$, then $d$ commutes with both $(cbab)^2$ and
$(abcb)^2$.
\end{lemma}
\begin{proof}
Put $x=bab$, $A=(ab)^2=ax$, and $F=(cd)^2$. As $a$ commutes with $F$,
the relation $[A,F]=1$ gives $[x,F]=1$. Also $[x,d]=1$, so
$[x,Fd]=[x,cdc]=1$. Conjugating by $c$ gives $[cxc,d]=1$.
Both factors of $(cbab)^2=(cxc)x$ commute with $d$. Finally
$abcb=b(cbab)^{-1}b$, and $b$ commutes with $d$.
\end{proof}

\begin{corollary}\label{cor:path}
The marked homomorphism $P\to H$ is an isomorphism. Furthermore $P_0$ is
the universal regular polytope group with the indicated consecutive square
sections, and $P$ is precisely its quotient by \eqref{eq:four}. Thus the
presentation assertion in [CFHS25, Theorem5.3] follows from
Theorem~\ref{thm:presentation}.
\end{corollary}
\begin{proof}
The complete local relations hold in $H$ by local recognition. Relations
\eqref{eq:four} hold in its base, which is class two. In each triple coordinate
one of the two disjoint edges loses an endpoint, and the square of its product
is the identity. Thus all relations hold in every coordinate, giving $P\to H$.
The local $E_i$ also map into $P$ by its defining relations.

Let $K_i$ be the image of $\gamma_3(E_i)$. It is normalized by the three
local markers. The two doubled translations are the words in
Lemma~\ref{lem:four}; that lemma and its reversed version show that the
nearest outside markers centralize them. All more distant markers commute
with every local marker by the string relations. Hence $K_i\unlhd P$.

After killing the $K_i$, each consecutive triple is class two. Every adjacent
square therefore commutes with its endpoints and its nearest outside markers;
it commutes with all remaining markers by the string relations. Thus the
quotient is class two, with the defining involutions, distant commutations and
adjacent product orders dividing four. It is a quotient of the complete
class-two path presentation of $B$. Conversely the map onward to the base
$B$ is a marked surjection. These two maps are inverse on generators, so this
quotient is exactly $B$. This proves the last condition of
Theorem~\ref{thm:presentation} without an upper bound on a finite order.
The theorem now identifies $P=H$, with \eqref{eq:pathorder} following from
the already established compatible-product formula.

For universality, every defining relator of $P_0$ survives all deletion
assignments. This is immediate for even Coxeter relators and is the proved
lattice-relator calculation in UA.Square for local relations. Hence $P_0$
has Boolean deletions and the intersection property by UA.D, without an
assumption of finiteness. Its local square subgroups embed because their maps
onward to $H$ are faithful; its markers remain nonidentity involutions.
It is therefore a string C-group. Every marked C-group with the prescribed
local sections satisfies its defining relations, and is its marked quotient.
The retained regular-polytope correspondence gives precisely the universal
polytope assertion. No order of $P_0$ is asserted.
\end{proof}

\section{What has been shared and what remains local}
The two applications used to prove separate upper bounds on a candidate
presentation and compare them to a constructed finite quotient. Their shared
replacement is now Theorem~\ref{thm:presentation}: it proves faithfulness
before any order is calculated, and also works for infinite local groups.
The underlying idea grows from the already reviewed UA.SK and UA.Order
arguments. We make no historical priority claim for the resulting abstract
group lemma. The geometric results retain their original attribution.

The genuine source-specific work is visible: in the hypertope case it is
\eqref{eq:parity}, the three normality relations and the square quotient; in
the string case it is Lemma~\ref{lem:four} and the class-two quotient
identification. Neither is removed by renaming it an umbrella input.
The supported-kernel construction, affine square models, and appropriate
geometric correspondence are still required dependencies. This note does not
replace the bounded cyclic-core classification in L035, the finite census and
presentation checks in N009, or all of F016. No source or theorem credit and no
compression ratio is assigned.

The Lean file \texttt{NormalKernelRecognition.lean} proves a finite-join
coordinate lemma, the joint injectivity assertion in
Theorem~\ref{thm:recognition}, and Corollary~\ref{cor:quotient}, in arbitrary
groups. It derives the required product decomposition from subgroup normality;
it does not take final injectivity as a premise. The companion file
\texttt{NormalKernelProduct.lean} constructs the explicit multiplicative
isomorphism $\prod_i K_i\simeq\ker\pi$. Its map is multiplication in $P$,
and on each individual factor it is the given subgroup inclusion. This last
conclusion uses equality in \eqref{eq:basekernel}; containment alone would
not imply it. Fourteen declaration reports use only the standard Lean axioms.
The formal construction uses mathlib's existing multiplication map for
commuting independent subgroups, also used in earlier project components.

These files do not formalize the source presentations, the construction of
UA.SK, the normality checks in the applications, the affine calculations or
the geometric correspondence in this note. Those remain the written
obligations displayed above. An existing Lean lemma,
\texttt{BooleanRetractions.four\_marker\_centralization}, verifies
Lemma~\ref{lem:four}; applying it to all local kernels still requires the
written argument in the final application.

\begin{thebibliography}{9}
\bibitem{UA}[UA] Human Compression Project, \emph{Supported extensions and finite
two-group quotients}, current Topic A umbrella; reviewed phase207 proofs,
canonical TeX edition. Required results are UA.D, UA.SK, its local, deletion
and geometric clauses, UA.Forest, UA.Affine, UA.Square and
UA.SquareDecomposition.
\bibitem{HFL}[HFL19] D.-D. Hou, Y.-Q. Feng and D. Leemans,
\emph{Existence of regular 3-hypertopes with $2^n$ chambers},
arXiv:1901.09034v1, 23 January 2019, Theorem3.2.
\bibitem{CFHS}[CFHS25] G. Cunningham, Y.-Q. Feng, D.-D. Hou and E. Schulte,
\emph{Symmetric polytopes whose automorphism groups are 2-groups},
arXiv:2512.15511v1, 18 December 2025, Theorems5.1 and5.3.
\bibitem{MS}[MS] P. McMullen and E. Schulte, \emph{Abstract Regular Polytopes},
Cambridge University Press, 2002, Section2E; the correspondence is retained
through the precise UA interface.
\end{thebibliography}
\end{document}
